\documentclass{article}
\usepackage{amsmath, amssymb, amsthm}
\usepackage{hyperref}
\usepackage{geometry}
\geometry{margin=1in}

\title{Erd\H{o}s Problem \#368}
\date{Last updated: 2025-08-31}
\author{Source: erdosproblems.com}

\begin{document}
\maketitle

\noindent\textbf{Status:} OPEN \quad \textbf{No prize}

\noindent\textbf{Tags:} number theory

\medskip

\noindent\textbf{Problem Statement:}

How large is the largest prime factor of $n(n+1)$?




Let $F(n)$ be the prime in question. P\'{o}lya \cite{Po18} proved that $F(n)\to \infty$ as $n\to\infty$. Mahler \cite{Ma35} showed that $F(n)\gg \log\log n$. Schinzel \cite{Sc67b} observed that for infinitely many $n$ we have $F(n)\leq n^{O(1/\log\log\log n)}$. The truth is probably $F(n)\gg (\log n)^2$ for all $n$. Erd\H{o}s \cite{Er76d} conjectured that, for every $\epsilon&gt;0$, there are infinitely many $n$ such that $F(n) &lt;(\log n)^{2+\epsilon}$. Pasten \cite{Pa24b} has proved that\[F(n) \gg \frac{(\log\log n)^2}{\log\log\log n}.\]The largest prime factors of $n(n+1)$ are listed as A074399 in the OEIS.




References

[Er76d] Erd\H{o}s, P., Problems and results on number theoretic properties of consecutive integers and related questions . Proceedings of the Fifth Manitoba Conference on Numerical Mathematics (Univ. Manitoba, Winnipeg, Man., 1975) (1976), 25-44.

[Ma35] Mahler, Kurt, \"{U}ber den gr\"{o}ssten Primteiler spezieller Polynome zweiten Grades . Archiv f\"{u}r math. og naturvid (1935).

[Pa24b] Pasten, Hector, The largest prime factor of {$n^2+1$} and improvements on subexponential {$ABC$} . Invent. Math. (2024), 373--385.

[Po18] P\'{o}lya, Georg, Zur arithmetischen {U}ntersuchung der {P}olynome . Math. Z. (1918), 143--148.

[Sc67b] Schinzel, A., On two theorems of Gelfond and some of their applications . Acta Arith. (1967/68), 177-236.

\noindent\textbf{Related OEIS sequences:} A074399


\bigskip
\noindent\small{Source: \url{https://www.erdosproblems.com/368}}
\end{document}
